% Part: proof-theory
% Chapter: sequent-calculus
% Section: introduction

\documentclass[../../../include/open-logic-section]{subfiles}

\begin{document}

\olfileid{pt}{seq}{int}

\olsection{પરિચય}

સિક્વન્ટ કલનો !!{proof} તંત્રો છે,
જેને ગેરહાર્ડ ગેન્ટ્ઝેને ૧૯૩૦ના દાયકામાં રજૂ કર્યાં.
તેમનો હેતુ વસ્તુઓને !!{proving} કરવાની વ્યવહારુ રીત આપવાનો ઓછો હતો.
ગેન્ટ્ઝેને તેમનો ઉપયોગ !!{proof}sની શક્ય સંરચના
અને ગુણધર્મો વિશે અમુક પરિણામો સાબિત કરી શકવા માટે કર્યો.
સાબિતી-સિદ્ધાંત બરાબર આવા પ્રશ્નોનો વિચાર કરે છે.

નામ સૂચવે છે તેમ સિક્વન્ટ કલનો \emph{સિક્વન્ટો} પર કામ કરે છે,
!!{formula}s પર નહીં.
સિક્વન્ટ એ !!{formula}sના બે ક્રમો કે બહુગણોની જોડી છે.
ગેન્ટ્ઝેનનાં મૂળ તંત્રો (\Log{LK} અને \Log{LJ})માં ક્રમો હતા;
હવે સાબિતી-સિદ્ધાંતના અભ્યાસીઓ બહુગણો પસંદ કરે છે.
બહુગણ સમૂહ જેવો છે, પરંતુ તેમાં !!{element}
એક કરતાં વધુ વાર આવી શકે છે.
સમૂહની જેમ તેમાં પણ !!{element}sનો ક્રમ મહત્ત્વનો નથી.
એટલે $\{!A, !A, !B\}$ અને $\{!A, !B, !A\}$
એ જ બહુગણ દર્શાવે છે,
પરંતુ તે બહુગણ $\{!A, !B\}$થી અને $\{!A, !A, !B, !B\}$થી જુદો છે,
કારણ કે છેલ્લાં બેમાં $!A$ અને $!B$ની સંખ્યા પહેલા બહુગણથી જુદી છે.

સાબિતી-સિદ્ધાંતમાં પરંપરાગત રીતે !!{formula}sના બહુગણો
મોટા ગ્રીક અક્ષરોથી લખાય છે.
એટલે $\Gamma$, $\Delta$, $\Pi$, $\Lambda$ કે $\Theta$ લખીએ,
તો આપણે જે તર્કનો વિચાર કરતાં હોઈએ તેનાં !!{formula}sના બહુગણો સમજીએ.
બીજી પરંપરા મુજબ $\tuple{\Gamma, \Delta}$ લખવાને બદલે
સાબિતી-સિદ્ધાંતના અભ્યાસીઓ બંને ઘટકોને સિક્વન્ટ-ચિહ્નથી જુદા પાડે છે;
અહીં~$\Sequent$ વાપરીએ છીએ.

\begin{defn}[સિક્વન્ટ]
\emph{સિક્વન્ટ} આ સ્વરૂપનો વ્યંજક છે:
\[
\Gamma \Sequent \Delta
\]
અહીં $\Gamma$ અને $\Delta$ એ !!{formula}sના સીમિત
(સંભવતઃ ખાલી) બહુગણો છે.
$\Gamma$ને \emph{પૂર્વાંગ} અને $\Delta$ને \emph{ઉત્તરાંગ} કહે છે.
\end{defn}

$\Gamma$નાં !!{formula}sનું સંયોજન $!G$~લઈએ
(ખાલી સંયોજન~$\ltrue$ ગણીએ)
અને $\Delta$નાં !!{formula}sનો વિકલ્પ $!D$~લઈએ
(ખાલી વિકલ્પ~$\lfalse$ ગણીએ).
ત્યારે સિક્વન્ટ $\Gamma \Sequent \Delta$
કોઈ સંરચના~$\Struct{M}$માં સત્ય છે
જો અને ફક્ત જો $!G \lif !D$ સત્ય હોય.
શાસ્ત્રીય તર્કમાં આ એ જ છે કે
કાં તો $\Gamma$નાં !!{formula}sમાંથી કોઈ અસત્ય હોય
અથવા $\Delta$નાં !!{formula}sમાંથી કોઈ સત્ય હોય.

$\Gamma$ એ !!{formula}sનો બહુગણ હોય,
તો $\Gamma$માં $!A$ ઉમેરવાનું પરિણામ
$\Gamma, !A$ (અથવા $!A, \Gamma$)થી લખીએ.
$\Delta$ પણ !!{formula}sનો બહુગણ હોય,
તો $\Gamma, \Delta$ બંને બહુગણોનો બહુગણ-સંઘ છે:
$\Gamma$માં $!A$ની \emph{આવૃત્તિઓની સંખ્યા}~n હોય
(એટલે તેમાં~$!A$ની $n$ નકલ હોય)
અને $\Delta$માં આવૃત્તિઓની સંખ્યા~$m$ હોય,
તો $\Gamma, \Delta$માં $!A$ની આવૃત્તિઓની સંખ્યા~$n+m$ છે.
સિક્વન્ટની એક બાજુનો બહુગણ ખાલી હોય,
તો સામાન્ય રીતે જગ્યા ખાલી રાખીએ.
દાખલા તરીકે $\Sequent !A$ એ ખાલી પૂર્વાંગવાળો સિક્વન્ટ છે,
જેના ઉત્તરાંગમાં $!A$ની આવૃત્તિઓની સંખ્યા~$1$ છે.
\sourcecorrection{OLMD-307}{મૂળ અહીં બે ક્રમોના બહુગણ-સંઘ કહે છે. આ પરિચ્છેદમાં બંને વસ્તુઓ બહુગણો છે અને આવૃત્તિઓની સંખ્યા ઉમેરાય છે; તેથી બંને બહુગણોનો સંઘ કહ્યું છે.}

સિક્વન્ટ કલનમાં !!^a{proof} એ સિક્વન્ટોનું વૃક્ષ છે.
સૌથી ઉપરના (અથવા આરંભિક) સિક્વન્ટો
વિચારેલા તંત્રના સ્વયંસિદ્ધો છે.
કોઈ સિક્વન્ટ બીજા એક કે બે સિક્વન્ટોની નીચે હોય,
તો તંત્રના કોઈ અનુમાન-નિયમથી તેમનામાંથી યોગ્ય રીતે અનુસરવો જોઈએ.
પહેલા શાસ્ત્રીય તર્ક માટેનાં બે જુદાં સિક્વન્ટ તંત્રો,
$\Log{G1c}$ અને~$\Log{G3c}$, વિચારીએ.
તેમના સ્વયંસિદ્ધો અને નિયમો
\cref{pt:seq:int:tab:G1c,pt:seq:int:tab:G3c}માં છે.
\oliflabeldef{fol:seq::chap}{ (ગેન્ટ્ઝેનનું મૂળ તંત્ર \Log{LK}
\olref[fol][seq][]{chap}માં વર્ણવ્યું છે.)}{}

\subfile{rules-G1c}
\subfile{rules-G3c}

તંત્રોમાં સૂક્ષ્મ તફાવતો છે, અને તેથી તેમના
સાબિતી-સૈદ્ધાંતિક ગુણધર્મો જુદા છે.
\Log{G1c}માં $!A \Sequent !A$ સ્વયંસિદ્ધો છે;
બીજાં !!{formula}s આવવા દેવાતાં નથી.
\Log{G3c}માં $!C, \Gamma \Sequent \Delta, !C$
સ્વરૂપનાં સ્વયંસિદ્ધો છે, જેમાં $\Gamma$ અને~$\Delta$માં
મનસ્વી ``બાજુનાં'' !!{formula}s હોય,
પરંતુ બંને બાજુ આવતું !!{formula}~$!C$ \emph{આણ્વિક} હોવું જોઈએ.
\Log{G1c}માં $\LeftR{\land}$ તથા $\RightR{\lor}$ના
બબ્બે સ્વરૂપ છે, જ્યારે \Log{G3c}માં એકેક છે.
વળી \Log{G1c}માં વધારાના ``બંધારણીય નિયમો'' છે,
જેને ``સંકોચન'' (\LeftR{\Contraction}, \RightR{\Contraction})
અને ``શિથિલીકરણ'' (\LeftR{\Weakening}, \RightR{\Weakening}) કહે છે.

\Log{G1c} અને \Log{G3c} તંત્રો શાસ્ત્રીય તર્ક માટે યથાર્થ અને પૂર્ણ છે.
બીજા શબ્દોમાં, આ તંત્રોમાં !!{provable} દરેક સિક્વન્ટ
દરેક !!{structure}માં સત્ય છે,
અને દરેક !!{structure}માં સત્ય હોય તે દરેક સિક્વન્ટની !!a{proof} છે.
એટલે કોઈ સિક્વન્ટની !!a{proof} છે \emph{કે નહીં} તે પ્રશ્નનો
જવાબ તર્કની બીજી રીતોથી પણ આપી શકાય
(દાખલા તરીકે નિદર્શ-સિદ્ધાંતની રીતોથી).
બંને તંત્રો દરેક !!{structure}માં સત્ય હોય
બરાબર તે જ સિક્વન્ટોને !!{prove} કરે છે;
તેથી ખાસ કરીને તેઓ એ જ સિક્વન્ટોને સાબિત કરે છે.

પરંતુ સાબિતી-સિદ્ધાંત ફક્ત કંઈક સાબિત થઈ શકે છે
\emph{કે નહીં} તેનો જ નહીં, \emph{કેવી રીતે} તેનો પણ વિચાર કરે છે.
તેના અભ્યાસીઓ આવા પ્રશ્નો વિચારે છે:
``કોઈ ચોક્કસ નિયમ હંમેશાં ટાળી શકાય?'',
``તંત્રમાં આ નિયમ ઉમેરવાથી નવા સિક્વન્ટો સાબિત થઈ ન જાય તે શક્ય છે?'',
અથવા ``એક તંત્રની !!{proof}sને બીજાની !!{proof}sમાં ફેરવી શકાય?''
જવાબ ``હા'' હોય, તો એ હકીકતની સાબિતી
ઘણી વાર !!{proof}sની જટિલતા પર આગમનથી થાય,
અથવા સમકક્ષ રીતે !!{proof}sની સંરચના પર આગમનથી થાય.
આથી ફક્ત હકીકતની સાબિતી જ નહીં
(દાખલા તરીકે \Log{G1c} સિક્વન્ટ સાબિત કરે તો \Log{G3c} પણ એ જ સાબિત કરે),
પણ પ્રક્રિયા પણ મળે
(દાખલા તરીકે~\Log{G1c}ની !!{proof}sને~\Log{G3c}ની !!{proof}sમાં ફેરવવાની).

સાબિતી-સિદ્ધાંતનું કેન્દ્રીય પરિણામ \emph{કટ-લોપનું પ્રમેય} છે.
તે બતાવે છે કે સિક્વન્ટ તંત્રોમાં કટ-નિયમ
\[
\Axiom$\Gamma \fCenter \Delta, !A$
\Axiom$!A, \Pi \fCenter \Lambda$
\RightLabel{\Cut}
\BinaryInf$\Gamma, \Pi \fCenter \Delta, \Lambda$
\DisplayProof
\]
ઉમેરી શકાય અને કયા સિક્વન્ટો !!{provable} છે તે બદલાતું નથી.
વળી તેની સાબિતી એવી રીત આપે છે
જે \Log{G1c} (અથવા \Log{G3c})ના નિયમો સાથે
\Cut{} વાપરતી કોઈ પણ !!{proof}ને,
ફક્ત \Log{G1c} (અથવા \Log{G3c})ના નિયમો વાપરતી સાબિતીમાં ફેરવે છે.
બીજા શબ્દોમાં આ પ્રક્રિયા \Cut{} વાપરતી !!{proof}sમાંથી
એ નિયમનો ``લોપ'' કરે છે; તેથી આ નામ છે.

\end{document}
